\setcounter{section}{10}

  \zwischenueberschrift{Soal latihan}

\inputaufgabe
{}
{

Tunjukkan bahwa

\mavergleichskettedisp
{\vergleichskette
{M
}
{ =} { \left\{ (x,y,z,t) \in \R^4  \mid   x+x^2y+z^2+t^3 = 0         \right\}
}
{ } {
}
{ } {
}
{ } {
}
}
{}{}{}
merupakan sebuah
\definitionsverweis {submanifold tertutup}{}{}
dari $\R^4$.

}
{} {}

\inputaufgabe
{}
{

Misalkan
\maabbdisp {f} {\R} {\R_+
} {}
sebuah
\definitionsverweis {fungsi yang terdiferensialkan secara kontinu}{}{.}
Tunjukkan bahwa
\definitionsverweis {permukaan putar}{}{}
dari
\definitionsverweis {grafik}{}{}
$f$ merupakan sebuah
\definitionsverweis {submanifold tertutup}{}{}
dari $\R^{3}$.

}
{} {}

\inputaufgabe
{}
{

Tunjukkan bahwa himpunan semua
\mathl{n\times n}{} matriks riil dengan
\definitionsverweis {determinan}{}{}
$1$ memiliki dimensi
\mathl{(n^2-1)}{} dan merupakan sebuah
\definitionsverweis {submanifold tertutup}{}{}
dari $\R^{n^2}$.

}
{} {}

\inputaufgabe
{}
{

Berikan sebuah contoh
\definitionsverweis {himpunan bagian tertutup}{}{}
\mathl{M\subseteq \R}{,} yang bukan
\definitionsverweis {submanifold tertutup}{}{} dari $\R$.

}
{} {}

\inputaufgabe
{}
{

Tentukan semua
\definitionsverweis {submanifold tertutup}{}{}
dari $\R$.

}
{} {}

\inputaufgabe
{}
{

Tentukan semua
\definitionsverweis {submanifold tertutup}{}{}
dari $S^1$.

}
{} {}

\inputaufgabe
{}
{

Misalkan
\mathkor {} {M} {dan} {N} {}
dua himpunan yang
\definitionsverweis {saling lepas}{}{}
dan masing-masing merupakan
\definitionsverweis {submanifold tertutup}{}{}
dari $\R^n$ dengan dimensi yang sama. Tunjukkan bahwa
\definitionsverweis {gabungan}{}{}
keduanya,

\mathl{M \cup N}{}, juga merupakan sebuah submanifold tertutup, dan bahwa pernyataan ini tidak berlaku tanpa asumsi saling lepas.

}
{} {}

\inputaufgabe
{}
{

Misalkan
\maabb {f} {\R^n } {\R
} {}
sebuah
\definitionsverweis {fungsi yang terdiferensialkan secara kontinu}{}{.}
Tunjukkan bahwa
\definitionsverweis {grafik}{}{}

\mavergleichskette
{\vergleichskette
{ \Gamma_f
}
{ \subseteq }{ \R^{n+1}
}
{  }{
}
{    }{
}
{   }{
}
}
{}{}{}
merupakan sebuah
\definitionsverweis {submanifold tertutup}{}{}
dari $\R^{n+1}$.

}
{} {}

\inputaufgabegibtloesung
{}
{

Misalkan

\mavergleichskette
{\vergleichskette
{ M
}
{ \neq }{ \emptyset
}
{  }{
}
{    }{
}
{   }{
}
}
{}{}{}
sebuah
\definitionsverweis {manifold diferensiabel}{}{}
berdimensi $n$. Tunjukkan bahwa terdapat sebuah rantai
\definitionsverweis {submanifold tertutup}{}{}

\mavergleichskettedisp
{\vergleichskette
{ M_0
}
{ \subseteq} { M_1
}
{ \subseteq} { M_2
}
{ \subseteq  \ldots \subseteq} { M_{n-1}
}
{ \subseteq} { M_n
}
}
{
\vergleichskettefortsetzung
{ =} { M
}
{ } {}
{ } {}
{ } {}
}{}{}
sedemikian sehingga submanifold tertutup $M_i$ berdimensi $i$.

}
{} {}

\inputaufgabegibtloesung
{}
{

Perhatikan himpunan

\mavergleichskettedisp
{\vergleichskette
{N
}
{ =} { \left\{ A= \begin{pmatrix}  a   &   b   \\  c  & d \end{pmatrix}   \mid   A \text{ adalah nilpoten}         \right\}
}
{ \subseteq} {\R^4
}
{ } {
}
{ } {
}
}
{}{}{}
yang terdiri atas
$2\times 2$ \definitionsverweis {matriks}{}{}
riil nilpoten, serta himpunan

\mavergleichskettedisp
{\vergleichskette
{M
}
{ =} { N \setminus \{ \begin{pmatrix}  0   &   0  \\  0  &  0 \end{pmatrix}  \}
}
{ } {
}
{ } {
}
{ } {
}
}
{}{}{.}

a) Apakah $N$ terhubung?

b) Tunjukkan bahwa $M$ merupakan sebuah
\definitionsverweis {submanifold tertutup}{}{}
dari suatu himpunan bagian terbuka

\mavergleichskette
{\vergleichskette
{ G
}
{ \subseteq }{ \R^4
}
{  }{
}
{    }{
}
{   }{
}
}
{}{}{.}

c) Tentukan
\definitionsverweis {dimensi}{}{}
$M$.

d) Apakah $M$ terhubung?

e) Tutupi
\mathl{M}{} dengan peta-peta topologis eksplisit.

}
{} {}

\inputaufgabe
{}
{

Misalkan

\mavergleichskette
{\vergleichskette
{ M
}
{ \subseteq }{ \R^n
}
{  }{
}
{    }{
}
{   }{
}
}
{}{}{}
sebuah
\definitionsverweis {submanifold tertutup}{}{}
dan

\mavergleichskette
{\vergleichskette
{ P
}
{ \in }{ M
}
{  }{
}
{    }{
}
{   }{
}
}
{}{}{}
sebuah titik. Misalkan
\mathl{T_P(i)}{} adalah
\definitionsverweis {pemetaan singgung}{}{}
yang bersesuaian dengan inklusi
\maabbdisp {i} { M } {\R^n
} {}
dan $\gamma$ sebuah
\definitionsverweis {lintasan diferensiabel}{}{}
di $M$ dengan

\mavergleichskettedisp
{\vergleichskette
{ \gamma(0)
}
{ =} { P
}
{ } {
}
{ } {
}
{ } {
}
}
{}{}{.}
Tunjukkan bahwa, di dalam

\mavergleichskette
{\vergleichskette
{ \R^n
}
{ \cong }{ T_P \R^n
}
{  }{
}
{    }{
}
{   }{
}
}
{}{}{}, kesamaan

\mavergleichskettedisp
{\vergleichskette
{ T_P(i) ( [\gamma] )
}
{ =} { \left(  i \circ \gamma  \right) '(0)
}
{ } {
}
{ } {
}
{ } {
}
}
{}{}{}
berlaku.

}
{} {}

\inputaufgabe
{}
{

Misalkan

\mavergleichskette
{\vergleichskette
{ G
}
{ \subseteq }{ \R^n
}
{  }{
}
{    }{
}
{   }{
}
}
{}{}{}
\definitionsverweis {terbuka}{}{,}
\maabb {\varphi} { G } {\R^m
} {}
sebuah
\definitionsverweis {pemetaan diferensiabel}{}{}
dan $M$ adalah
\definitionsverweis {serat}{}{}
di atas

\mavergleichskette
{\vergleichskette
{ 0
}
{ \in }{ \R^m
}
{  }{
}
{    }{
}
{   }{
}
}
{}{}{.}
Andaikan
\definitionsverweis {diferensial total}{}{}
bersifat
\definitionsverweis {surjektif}{}{}
di setiap titik serat ini. Tunjukkan bahwa untuk

\mavergleichskette
{\vergleichskette
{ P
}
{ \in }{ M
}
{  }{
}
{    }{
}
{   }{
}
}
{}{}{}
\definitionsverweis {ruang singgung}{}{}
dalam pengertian
[[Differenzierbare Abbildung/R/Regulärer Punkt/Tangentialraum/An Faser/Definition|Definisi 53.7 (Analisis (Osnabrück 2021-2023))  ]] berimpit dengan
\definitionsverweis {ruang singgung}{}{}
dari
\definitionsverweis {manifold diferensiabel}{}{}
$M$ di titik $P$.

}
{} {}

\inputaufgabe
{}
{

Misalkan $M$ sebuah
\definitionsverweis {manifold diferensiabel}{}{}
dan
\maabbdisp {\pi} {TM} {M
} {}
adalah
\definitionsverweis {bundel tangen}{}{.}
Tunjukkan bahwa pemetaan proyeksi ini
\definitionsverweis {kontinu}{}{.}

}
{} {}

\inputaufgabe
{}
{

Misalkan $TM$ adalah
\definitionsverweis {bundel tangen}{}{}
dari suatu
\definitionsverweis {manifold diferensiabel}{}{}
$M$. Tunjukkan bahwa himpunan-himpunan
\mathl{(T(\alpha))^{-1} (V \times W)}{} untuk semua peta
\maabbdisp {\alpha} { U } { V
} {}
dengan
\mathl{V,W \subseteq \R^n}{} terbuka membentuk suatu
\definitionsverweis {basis topologi}{}{}
pada bundel tangen tersebut.

}
{} {}

\inputaufgabegibtloesung
{}
{

Misalkan

\mavergleichskette
{\vergleichskette
{ W
}
{ \subseteq }{ \R^n
}
{  }{
}
{    }{
}
{   }{
}
}
{}{}{}
\definitionsverweis {terbuka}{}{,}
\maabb {h} { W } { \R
} {}
sebuah
fungsi kelas $C^2$
dan

\mavergleichskette
{\vergleichskette
{ Y
}
{ = }{ h^{-1}(c)
}
{  }{
}
{    }{
}
{   }{
}
}
{}{}{}
adalah
\definitionsverweis {serat}{}{}
di atas

\mavergleichskette
{\vergleichskette
{ c
}
{ \in }{ \R
}
{  }{
}
{    }{
}
{   }{
}
}
{}{}{,}
dengan $h$
\definitionsverweis {reguler}{}{}
di setiap titik pada $Y$. Realisasikan
\definitionsverweis {bundel tangen}{}{}
dari $Y$ sebagai sebuah
\definitionsverweis {submanifold tertutup}{}{}
dari $W \times \R^n$.

}
{} {}

\inputaufgabe
{}
{

Misalkan

\mavergleichskette
{\vergleichskette
{ Y
}
{ \subseteq }{ \R^n
}
{  }{
}
{    }{
}
{   }{
}
}
{}{}{}
sebuah
\definitionsverweis {submanifold tertutup}{}{}
dan
\maabb {\varphi} { V } { U \subseteq Y
} {}
sebuah parametrisasi difeomorfik dari suatu himpunan bagian terbuka

\mavergleichskette
{\vergleichskette
{ U
}
{ \subseteq }{ Y
}
{  }{
}
{    }{
}
{   }{
}
}
{}{}{}
oleh suatu himpunan bagian terbuka

\mavergleichskette
{\vergleichskette
{ V
}
{ \subseteq }{ \R^d
}
{  }{
}
{    }{
}
{   }{
}
}
{}{}{,}
dengan $\varphi$ juga dipandang sebagai pemetaan ke $\R^n$. Definisikan
\maabbdisp {X_i} {U} {TU
} {}
melalui $X_i(\varphi(x))=T_x(\varphi)(e_i)=\partial_i\varphi(x)$. Tunjukkan bahwa diagram

\mathdisp {\begin{matrix}  V  & \stackrel{ \partial_i }{\longrightarrow} &  TV &  \\ \!\!\!\!\! \,\, \, \varphi \downarrow & \partial_i \varphi \searrow & \downarrow T(\varphi) \!\!\!\!\!  & \\  U  & \stackrel{ X_i }{\longrightarrow} &  TU & \!\!\!\!\!   \\ \end{matrix}} {  }

tersebut komutatif.

}
{} {}

\inputaufgabe
{}
{

Misalkan
\mathkor {} {M} {dan} {N} {}
adalah
\definitionsverweis {manifold diferensiabel}{}{} dan
\maabb {\varphi} { M } { N
} {} sebuah
\definitionsverweis {pemetaan diferensiabel}{}{.} Tunjukkan bahwa
\definitionsverweis {pemetaan singgung}{}{}
yang bersesuaian,
\maabbdisp {T(\varphi)} {TM} {TN
} {}
\definitionsverweis {kontinu}{}{.}

}
{} {}

\inputaufgabe
{}
{

Hitung
\definitionsverweis {pemetaan singgung}{}{} $T\varphi$ untuk
\maabbeledisp {\varphi} {\R^3} {\R^2
} {(x,y,z)} {(x^2y-3xz^3+y^2,x \sin  y  -e^{yz})
} {} dengan menggunakan identifikasi
\mathkor {} {T\R^3 =\R^3 \times \R^3} {dan} {T\R^2 =\R^2 \times \R^2} {.}

}
{} {}

\inputaufgabe
{}
{

Berikan contoh sebuah kurva diferensiabel
\maabbdisp {\gamma} { [0,1[} { S^1
} {}
sedemikian sehingga
\definitionsverweis {limit}{}{}

\mathl{\operatorname{lim}_{   t  \rightarrow  1   } \,  \gamma(t)}{} ada, tetapi limit
\mathl{\operatorname{lim}_{   t  \rightarrow  1   } \, ( \gamma(t), T_t (\gamma ) (1))}{} di dalam
\mathl{TS^1}{} tidak ada.

}
{} {}

\inputaufgabe
{}
{

Misalkan

\mavergleichskette
{\vergleichskette
{ M
}
{ \subseteq }{ \R^n
}
{  }{
}
{    }{
}
{   }{
}
}
{}{}{}
sebuah
\definitionsverweis {submanifold tertutup}{}{.}
Tafsirkan komposisi

\mathdisp {TM \longrightarrow T\R^n = \R^n \times \R^n \stackrel{+}{\longrightarrow} \R^n} { . }

}
{} {}

\inputaufgabe
{}
{

Tunjukkan bahwa pemetaan
\maabbeledisp {} {TS^1} { \R^2
} {( (a,b),t (-b,a))} { (a,b) + t (-b,a)
} {,}
memiliki dua titik prabayang untuk setiap titik
\mathl{(x,y) \in \R^2}{} di luar cakram satuan, satu titik prabayang pada lingkaran satuan, dan tidak memiliki titik prabayang di dalam cakram satuan terbuka. Tafsirkan hasil ini secara geometris.

}
{} {}

\inputaufgabe
{}
{

Misalkan

\mavergleichskette
{\vergleichskette
{ M
}
{ \subseteq }{ N
}
{  }{
}
{    }{
}
{   }{
}
}
{}{}{}
sebuah
\definitionsverweis {submanifold tertutup}{}{}
dari
\definitionsverweis {manifold diferensiabel}{}{}
$N$. Tunjukkan bahwa

\mavergleichskette
{\vergleichskette
{ TM
}
{ \subseteq }{ TN
}
{  }{
}
{    }{
}
{   }{
}
}
{}{}{}
merupakan sebuah himpunan bagian tertutup.

}
{} {}

\inputaufgabe
{}
{

Berikan sebuah contoh pemetaan yang
\definitionsverweis {injektif}{}{}
dan
\definitionsverweis {diferensiabel}{}{},
\maabbdisp {\varphi} { M } { N
} {}
antara dua
\definitionsverweis {manifold diferensiabel}{}{},
\mathkor {} {M} {dan} {N} {}
sedemikian sehingga
\definitionsverweis {pemetaan singgung}{}{}
yang bersesuaian,
\maabbdisp {T(\varphi)} {TM} {TN
} {},
tidak injektif.

}
{} {}

\inputaufgabe
{}
{

Berikan sebuah contoh pemetaan yang
\definitionsverweis {surjektif}{}{}
dan
\definitionsverweis {diferensiabel}{}{},
\maabbdisp {\varphi} { M } { N
} {}
antara dua
\definitionsverweis {manifold diferensiabel}{}{},
\mathkor {} {M} {dan} {N} {}
sedemikian sehingga
\definitionsverweis {pemetaan singgung}{}{}
yang bersesuaian,
\maabbdisp {T(\varphi)} {TM} {TN
} {},
tidak surjektif.

}
{} {}

\inputaufgabegibtloesung
{}
{

Berikan sebuah
\definitionsverweis {medan vektor}{}{}
yang
\definitionsverweis {kontinu}{}{}
pada $S^2$ dan hanya memiliki satu titik nol.

}
{} {}

\inputaufgabe
{}
{

\definitionsverweis {Permukaan bola satuan}{}{}

\mathl{S^2}{} berputar satu putaran penuh dalam satu detik
\zusatzklammer {berlawanan arah jarum jam jika dilihat dari kutub utara} {} {}
mengelilingi sumbu kutub. Kurva diferensiabel dan vektor singgung manakah pada suatu titik
\mathl{P \in S^2}{} yang bersesuaian dengan gerak ini? Deskripsikan
\definitionsverweis {medan vektor}{}{}
yang terkait,
\maabbdisp {} {S^2} {TS^2 \subseteq T\R^3 = \R^6
} {.}

}
{} {}

Misalkan $M$ sebuah
\definitionsverweis {manifold diferensiabel}{}{}
dengan
\definitionsverweis {bundel tangen}{}{}
\maabb {p} {TM} {M
} {}
dan misalkan
\maabbdisp {\psi} { Z } { M
} {}
sebuah
\definitionsverweis {pemetaan}{}{,}
dengan $Z$ menyatakan suatu himpunan. Sebuah pemetaan
\maabbdisp {F} { Z } { TM
} {}
disebut
\definitionswort {medan vektor sepanjang}{}
$\psi$, jika

\mavergleichskette
{\vergleichskette
{ p \circ F
}
{ = }{ \psi
}
{  }{
}
{    }{
}
{   }{
}
}
{}{}{}
berlaku.

  \zwischenueberschrift{Soal untuk dikumpulkan}

\inputaufgabe
{8 (3+3+2)}
{

Misalkan

\mavergleichskette
{\vergleichskette
{ m,n
}
{ \in }{\N_+
}
{  }{
}
{    }{
}
{   }{
}
}
{}{}{.}

a) Tunjukkan bahwa himpunan

\mavergleichskettedisp
{\vergleichskette
{ M
}
{ =} { \left\{ (x,y,u,v) \in \R^4  \mid  ux^m+vy^n = 1         \right\}
}
{ } {
}
{ } {
}
{ } {
}
}
{}{}{}
merupakan sebuah
\definitionsverweis {submanifold tertutup}{}{}
dari $\R^4$.

b) Tunjukkan bahwa
\definitionsverweis {pemetaan}{}{}
\maabbeledisp {\varphi} { M } {\R^2
} {(x,y,u,v)} {(x,y)
} {,}
\definitionsverweis {diferensiabel}{}{}
dan di setiap titik

\mavergleichskette
{\vergleichskette
{ P
}
{ \in }{ M
}
{  }{
}
{    }{
}
{   }{
}
}
{}{}{}
\definitionsverweis {reguler}{}{.}

c) Deskripsikan
\definitionsverweis {serat-serat}{}{}
dari $\varphi$.

}
{} {}

\inputaufgabe
{10 (2+3+5)}
{

Perhatikan
\definitionsverweis {pemetaan}{}{}
\maabbeledisp {\varphi} {\R} {\R^2
} { x } {(x^2,x^3)
} {.}

a) Tunjukkan bahwa pemetaan ini
\definitionsverweis {diferensiabel}{}{}
dan
\definitionsverweis {injektif}{}{.}

b) Tunjukkan bahwa $\varphi$ tidak
\definitionsverweis {reguler}{}{}
pada setidaknya satu titik.

c) Tunjukkan bahwa
\definitionsverweis {citra}{}{}
$\varphi$
\definitionsverweis {tertutup}{}{}
di $\R^2$, tetapi bukan
\definitionsverweis {submanifold tertutup}{}{}
dari $\R^2$.

}
{} {}

\inputaufgabe
{4}
{

Tunjukkan bahwa terdapat sebuah
\definitionsverweis {homeomorfisme}{}{}
antara
\definitionsverweis {bundel tangen}{}{}
\mathl{TS^1}{} dari
$1$-\definitionsverweis {sfera}{}{} $S^1$ dan
\definitionsverweis {produk}{}{} $S^1 \times \R$.

}
{} {}

\inputaufgabe
{5}
{

Misalkan $M$ sebuah
\definitionsverweis {manifold diferensiabel}{}{}
dan
\maabbdisp {\pi} {TM} {M
} {}
adalah
\definitionsverweis {bundel tangen}{}{.}
Tunjukkan bahwa
\mathl{TM}{} sendiri secara alami merupakan sebuah
\definitionsverweis {manifold topologis}{}{.}

}
{} {}

Dalam soal berikut digunakan konsep modul atas $R$
\zusatzklammer {yang merupakan generalisasi langsung dari konsep ruang vektor} {} {.}

Misalkan $R$ sebuah
\definitionsverweis {gelanggang komutatif}{}{}
dan

\mavergleichskette
{\vergleichskette
{ M
}
{ = }{ (M,+,0)
}
{  }{
}
{    }{
}
{   }{
}
}
{}{}{}.
Dengan notasi \stichwort {aditif} {}, struktur ini merupakan sebuah
\definitionsverweis {grup komutatif}{}{.}
$M$ disebut sebuah
\definitionswortpraemath {R}{ modul }{,}
jika ditetapkan suatu operasi
\maabbeledisp {} {R \times M } { M
} {(r,v)} { rv = r\cdot v
} {,}
\zusatzklammer {yang disebut \stichwort {perkalian skalar} {}} {} {}
dan operasi itu memenuhi aksioma-aksioma berikut
\zusatzklammer {di sini

\mavergleichskettek
{\vergleichskettek
{ r,s
}
{ \in }{ R
}
{  }{
}
{    }{
}
{   }{
}
}
{}{}{}
dan

\mavergleichskettek
{\vergleichskettek
{ u,v
}
{ \in }{ M
}
{  }{
}
{    }{
}
{   }{
}
}
{}{}{}
dipilih sebarang} {} {:}
\aufzaehlungvier{
\mavergleichskette
{\vergleichskette
{ r(su)
}
{ = }{ (rs) u
}
{  }{
}
{    }{
}
{   }{
}
}
{}{}{,}
}{
\mavergleichskette
{\vergleichskette
{ r(u+v)
}
{ = }{ (ru) + (rv)
}
{  }{
}
{    }{
}
{   }{
}
}
{}{}{,}
}{
\mavergleichskette
{\vergleichskette
{ (r+s)u
}
{ = }{ (ru)+ (su)
}
{  }{
}
{    }{
}
{   }{
}
}
{}{}{,}
}{
\mavergleichskette
{\vergleichskette
{ 1 u
}
{ = }{  u
}
{  }{
}
{    }{
}
{   }{
}
}
{}{}{.}
}

\inputaufgabe
{4}
{

Misalkan $M$ sebuah
\definitionsverweis {manifold diferensiabel}{}{,}
misalkan

\mavergleichskette
{\vergleichskette
{ R
}
{ = }{ C^1(M,\R)
}
{  }{
}
{    }{
}
{   }{
}
}
{}{}{}
adalah
\definitionsverweis {gelanggang}{}{}
dari
\definitionsverweis {fungsi-fungsi diferensiabel}{}{}
pada $M$, dan misalkan $F$ adalah himpunan semua
\definitionsverweis {medan vektor}{}{}
pada $M$.

a) Definisikan suatu penjumlahan pada $F$ sedemikian sehingga $F$ menjadi sebuah
\definitionsverweis {grup komutatif}{}{.}

b) Definisikan suatu perkalian skalar
\maabbeledisp {} {R \times F } { F
} {(f,s)} {fs
} {,}
sedemikian sehingga $F$ menjadi sebuah
$R$-\definitionsverweis {modul}{}{.}

}
{} {}

 [[Kategorie:Latexseite]]
